\begin{abstract}
We present DriftJax, a one-dimensional drift--diffusion--Poisson framework
for photovoltaic simulation and design sensitivity analysis. The numerical
method combines Scharfetter--Gummel fluxes with an analytically assembled
block-tridiagonal Jacobian and a direct Block-Thomas Newton step. For a mesh
of $N$ nodes, the forward linear solve requires $O(N)$ work and storage.
Sensitivities of converged steady-state solutions are evaluated using the
implicit function theorem and a custom reverse-mode rule in JAX. The
production adjoint uses pivoted dense LU, with $O(N^3)$ factorisation work
and $O(N^2)$ matrix storage per bias; checkpointing reduces retained
intermediates without changing this scaling. This separation of forward
and backward linear algebra makes the computational trade-off explicit.
The framework combines carrier-statistics, optical-generation, and
recombination models with a common device interface. Reported verification
studies compare analytical Jacobians, linear residuals, conservation
identities, and finite-difference sensitivities. Illustrative applications
address derivative-based sensitivity analysis, constrained band-gap/thickness
optimisation, and recovery of a synthetic current--voltage curve.
The results support the use of structured forward computation within a
differentiable design workflow, while also identifying the importance of
conditioning, objective normalisation, and convergence checks. The reported
timing ratios are specific to the matched configurations and have not been
isolated by controlled ablation; the available release records do not
constitute a fully reproducible benchmark archive.
\end{abstract}